\documentclass[a4paper,11pt,twoside]{article}
\usepackage[utf8]{inputenc}
\usepackage{xparse}
\usepackage[english,greek]{babel}
\usepackage{alphabeta}
\usepackage{nimbusserif}
\usepackage[T1]{fontenc}
\usepackage[outer=1.50cm, inner=2.00cm, top=2.00cm, bottom=2.00cm]{geometry}
\usepackage{multicol,longtable,multirow,hhline,enumitem,tikz,pgfplots,tkz-euclide,tkz-tab,capt-of,fontawesome5,gensymb,tabularray,fancyhdr,etoolbox,eurosym,xcolor-material,siunitx,eqparbox,microtype,tasks}
\usetikzlibrary{arrows.meta}
\usepackage[most]{tcolorbox}
\usetikzlibrary{tikzmark}
\let\myBbbk\Bbbk
\let\Bbbk\relax
\usepackage[curlybraces,amsbb,mtphrb]{mtpro2}
\usepackage[explicit]{titlesec}
\usepackage{soul}
\newcommand{\eng}[1]{\selectlanguage{english}#1\selectlanguage{greek}}
\usepackage{mathimatika}
%\usepackage[usenames,dvipsnames,cmyk,table,x11names]{xcolor}
\def\xrwma{black}
\definecolor{steelblue}{cmyk}{.7,.278,0,.294}
\definecolor{doc}{cmyk}{1,0.455,0,0.569}
\definecolor{orange}{HTML}{ff7300}

\newcommand{\ekthetesdeiktes}{\DeclareMathSizes{10.95}{10.95}{7}{5}
\DeclareMathSizes{6}{6}{3.8}{2.7}
\DeclareMathSizes{8}{8}{5.1}{3.6}
\DeclareMathSizes{9}{9}{5.8}{4.1}
\DeclareMathSizes{10}{10}{6.4}{4.5}
\DeclareMathSizes{12}{12}{7.7}{5.5}
\DeclareMathSizes{14.4}{14.4}{9.2}{6.5}
\DeclareMathSizes{17.28}{17.28}{11}{7.9}
\DeclareMathSizes{20.74}{20.74}{13.3}{9.4}
\DeclareMathSizes{24.88}{24.88}{16}{11.3}

\makeatletter
\newcommand{\subsup}{
\AtBeginDocument{
\check@mathfonts
\fontdimen16\textfont2=2.5pt
\fontdimen17\textfont2=2.5pt
\fontdimen14\textfont2=4.5pt
\fontdimen13\textfont2=4.5pt}
}
\makeatother}
\usepackage{wrapfig}
\newenvironment{WrapText1}[3][r]
{\wrapfigure[#2]{#1}{#3}}
{\endwrapfigure}

\newenvironment{WrapText2}[3][l]
{\wrapfigure[#2]{#1}{#3}}
{\endwrapfigure}
\newcommand{\wrapr}[6]{
\begin{minipage}{\linewidth}\mbox{}\\
\vspace{#1}
\begin{WrapText1}{#2}{#3}
\vspace{#4}#5\end{WrapText1}#6
\end{minipage}}

\newcommand{\wrapl}[6]{
\begin{minipage}{\linewidth}\mbox{}\\
\vspace{#1}
\begin{WrapText2}{#2}{#3}
\vspace{#4}#5\end{WrapText2}#6
\end{minipage}}
\usepackage{etoolbox,hhline,moreenum}
\usepackage{caption} 
\captionsetup[table]{skip=5pt}
\makeatletter
\newif\ifLT@nocaption
\preto\longtable{\LT@nocaptiontrue}
\appto\endlongtable{%
\ifLT@nocaption
\addtocounter{table}{\m@ne}%
\fi}
\preto\LT@caption{%
\noalign{\global\LT@nocaptionfalse}}
\makeatother

\newlist{alist}{enumerate}{1}
\setlist[alist]{label=\let\textdexiakeraia\relax\alph*.}
\UseTblrLibrary{counter}

\pgfmathdeclarefunction{gauss}{2}{%
\pgfmathparse{1/(#2*sqrt(2*pi))*exp(-((x-#1)^2)/(2*#2^2))}%
}
\pgfkeys{/pgfplots/aks_on/.style={axis lines=center,
xlabel style={at={(current axis.right of origin)},xshift=1.5ex,anchor=center},
ylabel style={at={(current axis.above origin)},yshift=1.5ex, anchor=center}}}
\pgfkeys{/pgfplots/grafikh parastash/.style={red!80!black,line width=.4mm,samples=200}}
\pgfkeys{/pgfplots/belh ar/.style={tick label style={font=\scriptsize},axis line style={-latex}}}
\tikzstyle{pl}=[line width=0.3mm]
\tikzstyle{plm}=[line width=0.4mm]

\newcommand{\kerkissans}[1]{{\fontfamily{maksf}\selectfont {#1}}}
\AtBeginDocument{\renewcommand{\textstigma}{\textsigma\texttau}}

\definecolor{titlecolor}{HTML}{cd0f00}

\newbox\TitleUnderlineTestBox
\newcommand*\TitleUnderline[1]
{%
\bgroup
\setbox\TitleUnderlineTestBox\hbox{\colorbox{titleblue}\strut}%
\setul{\dimexpr\dp\TitleUnderlineTestBox-.3ex\relax}{.3ex}%
\ul{#1}%
\egroup
}
\newcommand*\SectionNumberBox[1]
{%
\colorbox{red!80!black}
{%
\makebox[2em][c]
{%
\color{white}%
\strut
\csname the#1\endcsname
}%
}%
\TitleUnderline{\ \ \ }%
}
\titleformat{\section}[hang]
{\Large\fontfamily{maksf}\selectfont}%
{\fcolorbox{black}{white}{%
\raisebox{0pt}[13pt][3pt]{\makebox[80pt]{% height, width
\color{black}{\kerkissans{\textbf{\thesection ο Κεφάλαιο}}}}%
}}}%
{0pt}%
{\fcolorbox{black}{black}{\raisebox{0pt}[13pt][3pt]{\color{white}\ \textbf{#1}\ }}}

\titleformat{\subsection}[hang]
{\large\bfseries\fontfamily{maksf}\selectfont}%
{\colorbox{red!80!black}{%
\raisebox{0pt}[13pt][3pt]{\makebox[30pt]{% height, width
\color{white}{\kerkissans{\textbf{\thesubsection}}}}%
}}}%
{0pt}%
{
%\colorbox{black}{\raisebox{0pt}[13pt][3pt]{\color{white}\ 
\ \textbf{#1}\ }
%}}

\makeatletter
\@addtoreset{section}{part}
\makeatother

\titleformat{\part}[display]
{\normalfont\huge\filcenter\bfseries}{}{-30pt}{\Huge \textcolor{red!80!black}{ \kerkissans{ #1}}}
\titlespacing*{\part} 
{0pt}{0pt}{0pt}

\setlist[enumerate]{itemsep=0mm,label=\textcolor{\xrwma}{\textbf{\thesection.\arabic*}}}
\definecolor{bblue}{HTML}{4F81BD}
\definecolor{rred}{HTML}{C0504D}
\definecolor{ggreen}{HTML}{9BBB59}
\definecolor{ppurple}{HTML}{9F4C7C}


\pgfdeclarelayer{background}
\pgfdeclarelayer{foreground}
\pgfsetlayers{background,main,foreground}

\newcommand{\pie}[4][]{
\begin{scope}[#1]
\pgfmathsetmacro{\curA}{90}
\pgfmathsetmacro{\r}{1}
\def\c{(0,0)}
\node[pie title] at (90:1.3) {#2};
\foreach \v/\s/\l in{#3}{
\pgfmathsetmacro{\deltaA}{\v/#4*360}
\pgfmathsetmacro{\nextA}{\curA + \deltaA}
\pgfmathsetmacro{\midA}{(\curA+\nextA)/2}

\path[slice,\s] \c
-- +(\curA:\r)
arc (\curA:\nextA:\r)
-- cycle;
\pgfmathsetmacro{\d}{max((\deltaA * -(.5/50) + 1) , .5)}

\begin{pgfonlayer}{foreground}
\path \c -- node[pos=\d,pie values,values of \s]{$\l$} +(\midA:\r);
\end{pgfonlayer}

\global\let\curA\nextA
}
\end{scope}
}

\newcommand{\legend}[2][]{
\begin{scope}[#1]
\path
\foreach \n/\s in {#2}
{
++(0,-10pt) node[\s,legend box] {} +(5pt,0) node[legend label] {\n}
}
;
\end{scope}
}
\definecolor{a}{cmyk}{0,1,1,0.05}
\definecolor{b}{cmyk}{0,.8,.8,.15}
\definecolor{c}{cmyk}{0,.8,.8,.0}
\definecolor{d}{cmyk}{0,.7,.7,0}
\definecolor{e}{cmyk}{0,.5,.5,0}

\pgfplotsset{every axis/.append style={
x tick label style={/pgf/number format/.cd, 1000 sep={.}}}}

\fancyhf{}
\newcommand{\myleftmark}{\leftmark}
\renewcommand{\myleftmark}{{\large Τυπολόγιο}}
\renewcommand{\headrulewidth}{1pt}
\renewcommand{\sectionmark}[1]{\markboth{\large Κεφάλαιο \thesection\ -\ #1}{} }

\makeatletter% so we can use macros with @ in their names
\ifthenelse{\boolean{@twoside}}{%
%\fancyhead[LE,RO]{%
%\begin{tikzpicture}[overlay, remember picture]%
%    \fill[\xrwma] (current page.north west) -- (current page.north)-- ($(current page.north)+(9mm,-.5in)$)-- ($(current page.north west)+(0,-.5in)$) -- cycle;
%    \node[anchor=north west, text=white, font=\Large\scshape, minimum size=1in, inner xsep=5mm] at ($(current page.north west)+(10mm,5mm)$) {\kerkissans{\textbf{\myleftmark}}};
%\fill[black] (current page.north east) -- (current page.north)-- ($(current page.north)+(9mm,-.5in)$)-- ($(current page.north east)+(0,-.5in)$) -- cycle;
%\node[anchor=north east, text=white, font=\scshape, minimum size=1in, inner xsep=5mm] at ($(current page.north east)+(-10mm,5mm)$) {\kerkissans{\textbf{\leftmark}}};
%\end{tikzpicture}
%}
\fancyhead[LE]{
\begin{tikzpicture}[overlay, remember picture]%
%    \fill[black] (current page.north west) -- (current page.north)-- ($(current page.north)+(9mm,-.5in)$)-- ($(current page.north west)+(0,-.5in)$) -- cycle;
    \node[anchor=north west, text=red4, font=\large\scshape, minimum size=1in, inner xsep=5mm] at ($(current page.north west)+(10mm,5mm)$) {\kerkissans{\textbf{\thepage\ $\lvert$ \leftmark}}};
%\fill[\xrwma] (current page.north east) -- (current page.north)-- ($(current page.north)+(9mm,-.5in)$)-- ($(current page.north east)+(0,-.5in)$) -- cycle;
\node[anchor=north east, text=black, font=\large\scshape, minimum size=1in, inner xsep=5mm] at ($(current page.north east)+(-15mm,5mm)$) {\kerkissans{\textbf{\myleftmark}}};
\end{tikzpicture}
}%
\fancyhead[RO]{
\begin{tikzpicture}[overlay, remember picture]%
%    \fill[\xrwma] (current page.north west) -- (current page.north)-- ($(current page.north)+(9mm,-.5in)$)-- ($(current page.north west)+(0,-.5in)$) -- cycle;
    \node[anchor=north west, text=red4, font=\large\scshape, minimum size=1in, inner xsep=5mm] at ($(current page.north west)+(15mm,5mm)$) {\kerkissans{\textbf{\myleftmark}}};
%\fill[black] (current page.north east) -- (current page.north)-- ($(current page.north)+(9mm,-.5in)$)-- ($(current page.north east)+(0,-.5in)$) -- cycle;
\node[anchor=north east, text=black, font=\scshape, minimum size=1in, inner xsep=5mm] at ($(current page.north east)+(-10mm,5mm)$) {\kerkissans{\textbf{\leftmark\ $ \lvert $\ \thepage}}};
\end{tikzpicture}
}%
}
\makeatother

%-------- ΠΑΡΑΤΗΡΗΣΕΙΣ -----------------
\newcounter{parathrhsh}[section]
\renewcommand{\theparathrhsh}{\arabic{parathrhsh}}  
\newcommand{\Parathrhsh}[1]{\refstepcounter{parathrhsh}{\textbf{\textcolor{white}{\faLightbulb}\ \ -\ \ \kerkissans{Παρατήρηση\hspace{2mm}\theparathrhsh}}}}{}

\newcommand{\tss}[1]{\textsuperscript{#1}}
\newcommand{\tssL}[1]{\MakeLowercase{\textsuperscript{#1}}}

%----------- ΠΑΡΑΤΗΡΗΣΗ------------------
\newenvironment{parat}[1]
{\begin{tcolorbox}[toptitle=1mm,
bottomtitle=1mm,title=\Parathrhsh,
breakable,
enhanced standard,lifted shadow={1mm}{-2mm}{3mm}{0.3mm}%
{black!50!white},
colback=red!5!white,
boxrule=0.1pt,
colframe=red!80!black,
fonttitle=\bfseries,width=#1]}
{\end{tcolorbox}}
%-----------------------------------------

%----------- ΑΣΚΗΣΗ ------------------
\newcommand\themanum{A}

\newcounter{thema}[section]
\renewcommand{\thethema}{\arabic{thema}}   
\newcommand{\Thema}{\refstepcounter{thema}{\textbf{\textcolor{black}{\faPenSquare\ \  \large \kerkissans{\textbf{\thethemaο\hspace{2mm}Θέμα}}}}}\hspace{0mm}}{}

\newlist{erwthma}{enumerate}{3}
\setlist[erwthma]{label=\bf{\kerkissans{\large{\textcolor{black}{\themanum.\arabic*}}}},itemsep=0mm,leftmargin=0.8cm}
%------------------------------------
\newenvironment{thema}[1]
{\begin{tcolorbox}[title=\Thema\kerkissans{\textbf{\large{  #1}}},breakable,
enhanced standard,titlerule=0pt,toprule=0pt, rightrule=0pt, bottomrule=0pt,
colback=white,left=0mm,top=1mm,bottom=0mm,
boxrule=0pt,
colframe=white,leftrule=0mm,sharp corners,coltitle=black]
\renewcommand{\themanum}{#1}}
{\end{tcolorbox}}
\newcommand{\dint}{\displaystyle\int}
\usepgfplotslibrary{fillbetween}
%-----------------------------------------

%----------- ΑΣΚΗΣΗ ------------------
\newcounter{askhsh}[section]
\renewcommand{\theaskhsh}{\thesection.\arabic{askhsh}}   
\newcommand{\Askhsh}{\refstepcounter{askhsh}{\textbf{\textcolor{\xrwma}{\faPenSquare\ \  \large \kerkissans{\textbf{Άσκηση\hspace{2mm}\theaskhsh}}}}}\hspace{1mm}}{}
%------------------------------------
\newenvironment{askhsh}[1]
{\begin{tcolorbox}[title=\Askhsh\ \ :\ \  {\textcolor{black}{\kerkissans{\bmath{#1}}}},breakable,
enhanced standard,titlerule=-.2pt,toprule=0pt, rightrule=0pt, bottomrule=0pt,
colback=white,left=2mm,top=1mm,bottom=0mm,
boxrule=0pt,
colframe=white,borderline west={1.5mm}{0pt}{\xrwma},leftrule=2mm,sharp corners,coltitle=\xrwma]}
{\end{tcolorbox}}
%-----------------------------------------

%------- ΣΤΥΛ ΠΑΡΑΔΕΙΓΜΑΤΟΣ -------
\newcounter{paradeigma}[section]
\renewcommand{\theparadeigma}{\kerkissans{\arabic{paradeigma}}}   
\newcommand{\Paradeigma}[1]{\refstepcounter{paradeigma}\kerkissans{\bmath{\textcolor{red!80!black}{\faPlay\large \ \ Παράδειγμα\hspace{2mm}\theparadeigma\;:\;}\hspace{1mm}  #1}}\\}{}
%-----------------------------------

%------- ΣΤΥΛ ΛΥΣΗΣ ------------------
\newcommand{\lysh}{\textcolor{\xrwma}{\kerkissans{\noindent\faCheck\ \textbf{ΛΥΣΗ}}}\\}
%------------------------------------

%-------- ΠΡΟΣΟΧΗ -----------------
\newcounter{prosoxi}[section]
\renewcommand{\theprosoxi}{\arabic{prosoxi}}  
\newcommand{\Prosoxi}[1]{\refstepcounter{prosoxi}{\faExclamationTriangle\ \ \ \textbf{Προσοχή\hspace{2mm}\thesection.\theprosoxi}}}{}

%----------- ΠΡΟΣΟΧΗ------------------
\newenvironment{prosoxi}[1]
{\begin{tcolorbox}[title=\Prosoxi,
breakable,
enhanced standard,lifted shadow={1mm}{-2mm}{3mm}{0.3mm}%
{black!50!white},
colback=red!5!white,
boxrule=0.1pt,
colframe=red!80!black,
fonttitle=\bfseries,width=#1]}
{\end{tcolorbox}}
%-----------------------------------------

\newcommand{\bhmata}{\textcolor{red!80!black}{\kerkissans{{\large \textbf{Βήματα}\\\vspace{-7mm}}}}}
\newlist{bhma}{enumerate}{3}
\setlist[bhma]{label=\bf{\textcolor{red!80!black}{\kerkissans{\arabic*\textsuperscript{o}\ :}}},leftmargin=0.9cm,itemindent=0cm,ref=\bf{\arabic*\textsuperscript{o}\;Βήμα}}

\newcommand{\tropoi}{\textcolor{red!80!black}{\kerkissans{{\large \textbf{Τρόποι}\\\vspace{-7mm}}}}}
\newlist{tropos}{enumerate}{3}
\setlist[tropos]{label=\bf{\textcolor{black}{\kerkissans{\arabic*\textsuperscript{oς} Τρόπος\ :}}},leftmargin=0cm,itemsep=0mm,itemindent=2.1cm,ref=\bf{\arabic*\textsuperscript{oς}\;Τρόπος}}
\newcommand{\en}[1]{{\selectlanguage{english}#1\selectlanguage{greek}}}
\newlist{periptwsh}{enumerate}{3}
\setlist[periptwsh]{label=\bf\textit{\arabic*\textsuperscript{oς}\;Περίπτωση :},leftmargin=0cm,itemsep=0mm,itemindent=2.8cm,ref=\bf{\arabic*\textsuperscript{oς}\;Περίπτωση}}
%---------- Θεώρημα --------------
\newcounter{Thewrhmabox}[section]
\renewcommand{\theThewrhmabox}{\arabic{Thewrhmabox}}
\newenvironment{Thewrhmabox}[2][\linewidth]
{\refstepcounter{Thewrhmabox}
\begin{tcolorbox}[breakable,
enhanced standard,
boxrule=0.7pt,titlerule=-.2pt,
width=\linewidth,
title style={color=white},
overlay unbroken and first={
\path[left color=\xrwma,right color=\xrwma!70!white]
([yshift=-\pgflinewidth]frame.north west) to ([yshift=-5pt]title.south west)[rounded corners=2pt] -- ([xshift=-#2-15pt,yshift=-5pt]title.south east) to[rounded corners=2pt] ([xshift=-#2,yshift=-\pgflinewidth]frame.north east) -- cycle;
},
fonttitle=\bfseries,
before=\par\medskip\noindent,
after=\par\medskip,
toptitle=3pt,
top=11pt,topsep at break=-5pt,
colback=white,title={\kerkissans{\faStop\ \ \large Θεώρημα \theThewrhmabox}} : {\textcolor{white}{\kerkissans{#1}}}]}
{\end{tcolorbox}}
%------------------------------------------

\newcounter{orismos}[section]
\renewcommand{\theorismos}{\arabic{orismos}}
\newcommand{\orism}{\refstepcounter{orismos}{\bf\kerkissans{\textcolor{\xrwma}{\faBook\ \ \large{Ορισμός}\hspace{2mm}\theorismos}}}\hspace{1mm}}{}

\newenvironment{Orismosbox}[1]
{\begin{tcolorbox}[title=\orism:\ \  {\bf{\large\kerkissans{#1}}},
breakable,
enhanced standard,
titlerule=-.2pt,
toprule=0pt, 
rightrule=0pt, 
bottomrule=0pt,
colback=white,
opacityfill=0,
left=2mm,
top=1mm,
bottom=0mm,
boxrule=0pt,
colframe=white,
borderline west={1.5mm}{0pt}{\xrwma},
leftrule=2mm,
sharp corners,
coltitle=black]}
{\end{tcolorbox}}
%------------------------------------------
\newcommand{\apanthsh}{{\kerkissans{\textbf{ΑΠΑΝΤΗΣΗ}}}\\}

\newcommand*\circled[1]{\tikz[baseline=(char.base)]{
            \node[shape=circle,draw,inner sep=2pt] (char) {#1};}}

\newcounter{idiothta}[section]
\renewcommand{\theidiothta}{\arabic{idiothta}}
\newcommand{\idiothta}[1]{\refstepcounter{idiothta}{\bf\kerkissans{\textcolor{\xrwma}{\faPlay\ \ \ \large{Ιδιότητα}\hspace{2mm}\theidiothta\ :\ }}}\hspace{1mm}{\bf\kerkissans{#1}}\\}{}

\newcounter{protashxa}[section]
\renewcommand{\theprotashxa}{\arabic{protashxa}}
\newcommand{\protashxa}[1]{\refstepcounter{protashxa}{\colorbox{black}{%
\raisebox{0pt}[10pt][3pt]{\makebox[70pt]{% height, width
\color{white}{\faEdit\ \kerkissans{Πρόταση \theprotashxa}}}}%
}\ \ \kerkissans{#1}\\[2mm]}}

\DeclareTblrTemplate{caption}{nocaptemplate}{}
\DeclareTblrTemplate{capcont}{nocaptemplate}{}
\DeclareTblrTemplate{contfoot}{nocaptemplate}{}
\NewTblrTheme{mytabletheme}{
  \SetTblrTemplate{caption}{nocaptemplate}{}
  \SetTblrTemplate{capcont}{nocaptemplate}{}
  \SetTblrTemplate{contfoot}{nocaptemplate}{}
}

\NewTblrEnviron{mytblr}
\SetTblrStyle{firsthead}{font=\bfseries}
\SetTblrStyle{firstfoot}{fg=red2}
\SetTblrOuter[mytblr]{theme=mytabletheme}
\SetTblrInner[mytblr]{
rowspec={t{7mm}},columns = {c},
  width = 0.85\linewidth,
  row{odd} = {bg=red9,fg=black,ht=8mm},
 row{even} = {bg=red7,fg=black,ht=8mm},
hlines={white},vlines={white},
row{1} = {bg=red4, fg=white, font=\bfseries\fontfamily{maksf}},rowhead = 1,
  hline{2} = {.7mm}, % midrule  
}

\DeclareRobustCommand{\officialeuro}{%
  \ifmmode\expandafter\text\fi
  {\fontencoding{U}\fontfamily{eurosym}\selectfont e}}

\titleformat{\paragraph}
{\normalfont\large}%
{}{0em}%
{{\color{black}\titlerule[0pt]}\vskip-.2\baselineskip{\parbox[t]{\dimexpr\textwidth-2\fboxsep\relax}{\raggedright\strut{{\textcolor{red!80!black}{\faSquare\ \ \kerkissans{\bmath{#1}}}}}\strut}}}[\vskip -.2\baselineskip{}]
\setlength{\parindent}{0pt}

\begin{document}
\begin{thema}{Δ}
Δίνεται η συνάρτηση $f(x) = (x-1)\ln x - x + 1$, με $x > 0$.
\begin{erwthma}
    \item Να μελετήσετε τη συνάρτηση $f$ ως προς τη μονοτονία και τα ακρότατα.
    \item Να αποδείξετε ότι για κάθε $x>0$ ισχύει $(x-1)\ln x \ge x - 1$. Πότε ισχύει η ισότητα;
    \item Να βρείτε την εξίσωση της εφαπτομένης της $C_f$ η οποία διέρχεται από την αρχή των αξόνων $O(0,0)$.
    \item Θεωρούμε την αρχική συνάρτηση $F$ της $f$ στο $(0, +\infty)$. Να αποδείξετε ότι για κάθε $\alpha, \beta > 1$ με $\alpha \ne \beta$ ισχύει:
    \[ \frac{F(\beta) - F(\alpha)}{\beta - \alpha} > 0 \]
\end{erwthma}
\end{thema}

%---------------------------------------------------------------
\begin{thema}{Δ}
Έστω η συνάρτηση $f(x) = e^x - \ln(x+1) - 1$, ορισμένη στο $(-1, +\infty)$.
\begin{erwthma}
    \item Να αποδείξετε ότι η συνάρτηση $f$ είναι αυστηρά κυρτή στο $(-1, +\infty)$.
    \item Να αποδείξετε ότι η εξίσωση $f'(x) = 0$ έχει ακριβώς μία ρίζα $x_0 \in \left(-\frac{1}{2}, 0\right)$.
    \item Να μελετήσετε την $f$ ως προς τη μονοτονία, να βρείτε το είδος του ακροτάτου της και να αποδείξετε ότι $f(x) \ge f(x_0)$ για κάθε $x > -1$.
    \item Αν $F$ είναι μια παράγουσα (αρχική) της $f$ στο $(-1, +\infty)$, να δείξετε ότι για κάθε $x > 0$ ισχύει: $F(3x) + F(x) > 2F(2x)$.
\end{erwthma}
\end{thema}

%---------------------------------------------------------------
\begin{thema}{Δ}
Δίνεται η συνεχής συνάρτηση $f: \mathbb{R} \to \mathbb{R}$ για την οποία ισχύει:
\[ f^3(x) + 3f(x) = x + 4, \quad \text{για κάθε } x \in \mathbb{R} \]
\begin{erwthma}
    \item Να αποδείξετε ότι η $f$ είναι συνάρτηση 1-1 και να βρείτε την αντίστροφή της $f^{-1}$.
    \item Να μελετήσετε την $f^{-1}$ ως προς τη μονοτονία και να βρείτε τα κοινά σημεία των γραφικών παραστάσεων $C_f$ και $C_{f^{-1}}$.
    \item Να υπολογίσετε το όριο $\displaystyle{\lim_{x\to +\infty} \frac{f(x)}{\sqrt[3]{x}}}$.
    \item Να υπολογίσετε το εμβαδόν του χωρίου που περικλείεται από τις $C_f, C_{f^{-1}}$.
\end{erwthma}
\end{thema}

%---------------------------------------------------------------
\begin{thema}{Δ}
Δίνεται η συνάρτηση $f(x) = \frac{x^2}{x^2+1}$ με $x \in \mathbb{R}$.
\begin{erwthma}
    \item Να βρείτε τις ασύμπτωτες της γραφικής παράστασης της συνάρτησης $f$.
    \item Να μελετήσετε την $f$ ως προς τα ακρότατα και να προσδιορίσετε τα σημεία καμπής της.
    \item Να κατασκευάσετε έναν πρόχειρο πίνακα μεταβολών και να σχεδιάσετε τη γραφική της παράσταση.
    \item Να υπολογίσετε το εμβαδόν του χωρίου που περικλείεται από τη γραφική παράσταση της $f$, τον άξονα $x'x$ και τις ευθείες $x=0, x=1$.
\end{erwthma}
\end{thema}

%---------------------------------------------------------------
\begin{thema}{Δ}
Δίνεται η συνάρτηση $f(x) = x e^x - e^x + 1$, με $x \in \mathbb{R}$.
\begin{erwthma}
    \item Να μελετήσετε την $f$ ως προς τη μονοτονία και να λύσετε την εξίσωση $f(x) = 0$.
    \item Να αποδείξετε ότι η $C_f$ δέχεται οριζόντια ασύμπτωτη στο $-\infty$.
    \item Να αποδείξετε ότι η εξίσωση $x e^x - e^x + 1 = 2026$ έχει ακριβώς μία ρίζα στο $\mathbb{R}$.
    \item Να υπολογίσετε το εμβαδόν του χωρίου που περικλείεται από τη γραφική παράσταση της $f$, τον άξονα $x'x$ και τις ευθείες $x=0, x=1$.
\end{erwthma}
\end{thema}

%---------------------------------------------------------------
\begin{thema}{Δ}
Έστω μια παραγωγίσιμη συνάρτηση $f: (0, +\infty) \to \mathbb{R}$ για την οποία ισχύει:
\[ x f'(x) + f(x) = x e^x, \quad \text{για κάθε } x > 0 \quad \text{και} \quad f(1) = e \]
\begin{erwthma}
    \item Να αποδείξετε ότι ο τύπος της συνάρτησης είναι $f(x) = \frac{x e^x - e^x + e}{x}$.
    \item Να αποδείξετε ότι η συνάρτηση $f$ είναι γνησίως αύξουσα.
    \item Να υπολογίσετε τα όρια $\displaystyle\lim_{x\to 0^+} f(x)$ και $\displaystyle\lim_{x\to +\infty} f(x)$.
    \item Να δείξετε ότι η εφαπτομένη της γραφικής παράστασης $C_f$ στο $x_0 = 1$ τέμνει τον άξονα $y'y$ στο σημείο με τεταγμένη ίση με το $\lim_{x\to 0^+} f(x)$.
\end{erwthma}
\end{thema}

%---------------------------------------------------------------
\begin{thema}{Δ}
Θεωρούμε τη συνάρτηση $f(x) = \ln\left( x + \sqrt{x^2+1} \right)$, με $x \in \mathbb{R}$.
\begin{erwthma}
    \item Να αποδείξετε ότι η $f$ είναι περιττή συνάρτηση και ότι η $C_f$ διέρχεται από την αρχή των αξόνων.
    \item Να μελετήσετε τη μονοτονία της $f$ και να δείξετε ότι $f'(x) = \frac{1}{\sqrt{x^2+1}}$. 
    \item Να συζητήσετε την κυρτότητα της $f$ και να αναδείξετε ένα μοναδικό σημείο καμπής.
    \item Να λύσετε την ανίσωση $f(x^2 - 4) + f(x - 2) < 0$.
    \item Να αποδείξετε ότι το εμβαδόν μεταξύ $C_f$ και άξονα $x'x$ στο $[-\alpha, \alpha]$ ($\alpha>0$) ισούται με $2 \dint_0^\alpha f(x) dx$.
\end{erwthma}
\end{thema}

%---------------------------------------------------------------
\begin{thema}{Δ}
Δίνεται η συνάρτηση $f(x) = e^{-x} \sin x$, με $x \in [0, \pi]$.
\begin{erwthma}
    \item Να βρείτε τις ρίζες της $f$ στο διάστημα $[0, \pi]$.
    \item Να μελετήσετε την $f$ ως προς τα τοπικά ακρότατα στο $[0, \pi]$ και να αποδείξετε ότι η μέγιστη τιμή της είναι $e^{-\pi/4} \frac{\sqrt{2}}{2}$.
    \item Να αποδείξετε ότι υπάρχει $\xi \in (0, \pi)$ τέτοιο ώστε η εφαπτομένη της $C_f$ στο $x = \xi$ να είναι παράλληλη στον άξονα $x'x$.
    \item Να υπολογίσετε το ορισμένο ολοκλήρωμα $\dint_0^\pi f(x) dx$ χρησιμοποιώντας διαδοχική ολοκλήρωση κατά παράγοντες.
\end{erwthma}
\end{thema}

%---------------------------------------------------------------
\begin{thema}{Δ}
Έστω $f: \mathbb{R} \to \mathbb{R}$ μια δύο φορές παραγωγίσιμη συνάρτηση για την οποία ισχύει: $f(0)=f(1)=0$ και $f''(x) > 0$ για κάθε $x \in \mathbb{R}$.
\begin{erwthma}
    \item Να αποδείξετε ότι $f(x) < 0$ για κάθε $x \in (0, 1)$.
    \item Να δείξετε ότι υπάρχει ακριβώς ένα $x_0 \in (0, 1)$ τέτοιο ώστε η εφαπτομένη της $C_f$ στο $x_0$ να είναι παράλληλη με τον άξονα $x'x$.
    \item Να αποδείξετε ότι η $f$ παρουσιάζει ένα μοναδικό ολικό ελάχιστο.
    \item Αν επιπλέον $f'(0)=-1$ και $f'(1)=2$, να υπολογίσετε το ολοκλήρωμα $\dint_0^1 x f''(x) dx$.
\end{erwthma}
\end{thema}

%---------------------------------------------------------------
\begin{thema}{Δ}
Δίνεται η συνάρτηση $f(x) = \frac{\ln x}{x}$, με $x > 0$.
\begin{erwthma}
    \item Να μελετήσετε την μονοτονία και να βρείτε το ολικό μέγιστο της $f$.
    \item Με χρήση του παραπάνω ακροτάτου, να αποδείξετε ότι $x^e \le e^x$ για κάθε $x > 0$.
    \item Να βρείτε τις ασύμπτωτες (οριζόντιες και κατακόρυφες) της $C_f$.
    \item Να συγκρίνετε τους αριθμούς $\pi^e$ και $e^\pi$. (Δικαιολογήστε αναλυτικά την απάντησή σας).
    \item Αν $\alpha > \beta > e$, τότε να αποδείξετε ότι $f(\alpha) < f(\beta)$ και συνεπώς να διατάξετε τα $\alpha^\beta$ και $\beta^\alpha$.
\end{erwthma}
\end{thema}

%---------------------------------------------------------------
\begin{thema}{Δ}
Θεωρούμε τη συνάρτηση $f(x) = x^3 - 3x + \alpha$, όπου $\alpha \in \mathbb{R}$.
\begin{erwthma}
    \item Να βρείτε τα τοπικά ακρότατα της $f$ με βάση τη μονοτονία της.
    \item Να προσδιορίσετε τις τιμές του πραγματικού αριθμού $\alpha$, ώστε η εξίσωση $f(x) = 0$ να έχει ακριβώς τρεις πραγματικές, άνισες μεταξύ τους ρίζες.
    \item Έστω ότι $\alpha=2$. Να δείξετε ότι η $C_f$ εφάπτεται του άξονα $x'x$ σε ένα σημείο και τέμνει τον άξονα σε ένα άλλο.
    \item Για $\alpha=2$, να υπολογίσετε το εμβαδόν του χωρίου που περικλείεται από την $C_f$ τον άξονα $y'y$ τον άξονα $x'x$ για τα $x>0$.
\end{erwthma}
\end{thema}

%---------------------------------------------------------------
\begin{thema}{Δ}
Έστω η συνάρτηση $f(x) = -x + \sqrt{x^2+1}$ για $x \in \mathbb{R}$.
\begin{erwthma}
    \item Να βρείτε την παράγωγο $f'$ και να αποδείξετε ότι $f'(x) < 0$, άρα η $f$ είναι γνησίως φθίνουσα στο $\mathbb{R}$.
    \item Να υπολογίσετε το $\lim_{x\to +\infty} f(x)$ και να δείξετε ότι η ευθεία $y=0$ είναι οριζόντια ασύμπτωτη.
    \item Να βρείτε το $\lim_{x\to -\infty} f(x)$ και να δείξετε ότι η ευθεία $y = -2x$ είναι πλάγια ασύμπτωτη της $C_f$ στο $-\infty$.
    \item Να αποδείξετε ότι $f(x) > 0$ για κάθε πραγματικό αριθμό $x$.
\end{erwthma}
\end{thema}

%---------------------------------------------------------------
\begin{thema}{Δ}
Έστω παραγωγίσιμη συνάρτηση $f: \mathbb{R} \to \mathbb{R}$ τέτοια ώστε:
\[ x f'(x) + f(x) + e^x = 0 \quad \text{για κάθε } x \ne 0 \quad \text{και} \quad f(1) = 1-e \]
\begin{erwthma}
    \item Να αποδείξετε ότι $(x f(x) + e^x)' = 0$ για κάθε $x \neq 0$ και να βρείτε τον τύπο της $f(x)$ για $x \ne 0$.
    \item Έστω $f(x) = \frac{1 - e^x}{x}$ για $x \ne 0$. Αφού η $f$ είναι συνεχής στο $\mathbb{R}$, να βρείτε την τιμή $f(0)$.
    \item Να δείξετε ότι η $f$ είναι γνησίως φθίνουσα.
    \item Να υπολογιστεί το όριο $\displaystyle \lim_{x \to -\infty} f(x)$  και να προσδιοριστεί η ασύμπτωτη.
\end{erwthma}
\end{thema}

%---------------------------------------------------------------
\begin{thema}{Δ}
Δίνεται η παραγωγίσιμη συνάρτηση $f:(0, +\infty)\to\mathbb{R}$ για την οποία ισχύει: $f'(x) = \frac{\ln x}{x^2}$ και $f(1)=0$.
\begin{erwthma}
    \item Να βρείτε τον τύπο της $f(x)$. \textit{(Απάντηση: $f(x) = -\frac{\ln x}{x} - \frac{1}{x} + 1$)}.
    \item Να μελετηθεί η οικεία συνάρτηση $f$ ως προς μονοτονία, να βρεθούν τα τοπικά ακρότατα και να υπολογιστεί η μέγιστη τιμή.
    \item Να δειχτεί ότι η εξίσωση $f(x) = \frac{1}{2}$ έχει ακριβώς δύο ρίζες στο $(0, +\infty)$.
    \item Να υπολογιστεί το εμβαδόν μεταξύ της καμπύλης, του άξονα $x'x$ και της ευθείας που διέρχεται από το μέγιστο της συνάρτησης, παράλληλα προς τον $y'y$.
\end{erwthma}
\end{thema}

%---------------------------------------------------------------
\begin{thema}{Δ}
Θεωρούμε τη συνάρτηση $f(x) = x^2e^{-x}$.
\begin{erwthma}
    \item Να μελετήσετε τη μονοτονία της διευκρινίζοντας τα τοπικά ακρότατα.
    \item Να βρείτε τα σημεία καμπής της $f$ και τις ασύμπτωτες της γραφικής παράστασης (αν υπάρχουν).
    \item Να σχεδιάσετε με ακρίβεια τη γραφική παράσταση (πίνακας μεταβολών).
    \item Να δείξετε ότι για κάθε $x \in \mathbb{R}$ ισχύει η ανισότητα $x^2 \le e^{x-2} \cdot 4$. 
    \item Το εμβαδόν $\Omega$ ανάμεσα στην $C_f$, τον άξονα $x'x$ και τις ευθείες $x=0, x=\alpha$ με $\alpha>0$. Να δείξετε ότι $\lim_{\alpha\to+\infty} E(\alpha) = 2$.
\end{erwthma}
\end{thema}

\begin{thema}{Δ}
Δίνεται η συνάρτηση $f(x) = \begin{cases} \frac{x \ln x}{x-1}, & 0 < x \ne 1 \\ 1, & x = 1 \end{cases}$
\begin{erwthma}
    \item Να αποδείξετε ότι η συνάρτηση $f$ είναι συνεχής στο σημείο $x_0 = 1$.
    \item Να αποδείξετε ότι η συνάρτηση $f$ είναι παραγωγίσιμη στο $x_0 = 1$ και να βρείτε την $f'(1)$.
    \item Να αποδείξετε ότι παρουσιάζει κατακόρυφη ασύμπτωτη στο $0$ και να καθορίσετε τη συμπεριφορά της.
    \item Να δείξετε ότι $f'(x) > 0$ για $x > 0$, $x \ne 1$, επομένως είναι γνησίως αύξουσα.
\end{erwthma}
\end{thema}

%---------------------------------------------------------------
\begin{thema}{Δ}
Έστω η δύο φορές παραγωγίσιμη συνάρτηση $f:\mathbb{R}\to\mathbb{R}$, για την οποία $f(0)=0$ και 
\[ f''(x) - 2f'(x) + f(x) = 0 \quad \text{και} \quad f'(0) = 1 \]
\begin{erwthma}
    \item Να αποδείξετε ότι η συνάρτηση $g(x) = f(x) e^{-x}$ ικανοποιεί $g''(x)=0$ για κάθε $x \in \mathbb{R}$.
    \item Να βρείτε τον τύπο της $f(x)$ καταλήγοντας στον $f(x) = x e^x$.
    \item Να μελετήσετε την $f$ ως προς μονοτονία, ακρότατα, κυρτότητα.
    \item Να λύσετε την ανίσωση $f(x^2) - f(3x) < 0$ στο $\mathbb{R}$.
\end{erwthma}
\end{thema}

%---------------------------------------------------------------
\begin{thema}{Δ}
Θεωρούμε την συνάρτηση $f(x) = (x-2)e^x + x + 2$, με $x \in \mathbb{R}$.
\begin{erwthma}
    \item Να εξετάσετε τη συνάρτηση ως προς τη μονοτονία, να βρείτε τις ρίζες της (αν υπάρχουν) και το πρόσημό της.
    \item Σε ποιο σημείο της $C_f$ η εφαπτομένη προσεγγίζει τη μικρότερη δυνατή κλίση;
    \item Να βρείτε την πλάγια ασύμπτωτη της γραφικής παράστασης $C_f$ στο $-\infty$.
    \item Να βρείτε το εμβαδόν του χωρίου που περικλείεται από την $C_f$, την ασύμπτωτη της στο $-\infty$, και τις ευθείες $x=0$, $x=-1$.
\end{erwthma}
\end{thema}

%---------------------------------------------------------------
\begin{thema}{Δ}
Έστω η συνάρτηση $f(x) = \ln(x^2+1) - x$, με $x \in \mathbb{R}$.
\begin{erwthma}
    \item Να αποδείξετε ότι η συνάρτηση $f$ είναι γνησίως φθίνουσα στο $\mathbb{R}$ και να λύσετε την εξίσωση $f(e^x - 1) = f(0)$.
    \item Να μελετήσετε την $f$ ως προς την κυρτότητα και σημεία καμπής.
    \item Να δείξετε ότι η $C_f$ βρίσκεται "κάτω" από την εφαπτομένη της στο σημείο καμπής.
    \item Υπολογίστε το $\lim_{x\to +\infty} \frac{f(x) + x}{x^2}$.
\end{erwthma}
\end{thema}

%---------------------------------------------------------------
\begin{thema}{Δ}
Δίνεται η συνάρτηση $f(x) = \frac{e^x}{x^2+1}$ με $x \in \mathbb{R}$.
\begin{erwthma}
    \item Να προσδιορίσετε την μονοτονία και τα ακρότατα της $f$. \textit{(Προσοχή: ρίζα $x=1$ στην παράγωγο)}.
    \item Να αποδείξετε ότι ισχύει $\frac{e^x}{x^2+1} \ge \frac{e}{2}$ για κάθε $x \in \mathbb{R}$.
    \item Να βρείτε το σύνολο τιμών της $f$. (Χρειάζεται εύρεση ασυμπτώτων $\pm\infty$).
    \item Να λύσετε την εξίσωση $f(\ln x) - f(\frac{1}{x}) = 0$.
\end{erwthma}
\end{thema}

%---------------------------------------------------------------
\begin{thema}{Δ}
Σε μία χημική αντίδραση το ποσό μίας ουσίας $Q(t)$ μειώνεται βάσει του νόμου $Q'(t) = -k Q(t)$ \textbf{ή} εναλλακτικά ας μελετήσουμε τη γεωμετρική διαφορική $f'(x) = f(x) e^x$ με $f(0)=e$.
\begin{erwthma}
    \item Να αποδείξετε ότι η $f(x) = e^{e^x}$ επαληθεύει την αρχική συνθήκη της διαφορικής.
    \item Να μελετήσετε την $f$ (κυρτότητα και μονοτονία).
    \item Να βρείτε το $\lim_{x\to -\infty} f(x)$ και να δείξετε ότι υπάρχει οριζόντια ασύμπτωτη.
    \item Να βρείτε την εξίσωση εφαπτομένης της $\ln f(x)$ στο $x=0$.
\end{erwthma}
\end{thema}

%---------------------------------------------------------------
\begin{thema}{Δ}
Θεωρούμε την συνάρτηση $f(x) = 2x + \ln(x^2+1)$.
\begin{erwthma}
    \item Να αποδείξετε ότι η συνάρτηση $f$ αντιστρέφεται και να βρείτε το πεδίο ορισμού της $f^{-1}$.
    \item Να λύσετε την ανίσωση $f(x^2-x) > f(x+3)$.
    \item Να εξετάσετε την $f$ ως προς την κυρτότητα και να βρείτε την εφαπτομένη της στο σημείο καμπής.
    \item Να υπολογίσετε το εμβαδόν μεταξύ της γραφικής παράστασης $C_f$, του άξονα $x'x$ και των ευθειών $x=0, x=1$.
\end{erwthma}
\end{thema}

%---------------------------------------------------------------
\begin{thema}{Δ}
Έστω η συνάρτηση $f(x) = \left(x + \frac{1}{x}\right) \ln x$, με $x > 0$.
\begin{erwthma}
    \item Να εξετάσετε αν η ευθεία $x=0$ αποτελεί κατακόρυφη ασύμπτωτη της $C_f$.
    \item Να βρείτε την $f'(x)$ και να αποδείξετε: $\alpha$. $x^2 -1 + 2\ln x > 0$ για $x>1$. $\beta$. Μονοτονία της $f$. 
    \item Να λύσετε την (κοινή) ανίσωση $\left(x + \frac{1}{x}\right) \ln x > 0$.
    \item Να υπολογίσετε το $\dint_1^e \frac{f(x)}{x} dx$.
\end{erwthma}
\end{thema}

%---------------------------------------------------------------
\begin{thema}{Δ}
Δίνεται η συνάρτηση $f(x) = e^x - ax$, με $\alpha>0$ σταθερό πραγματικό αριθμό.
\begin{erwthma}
    \item Να μελετήσετε τη μονοτονία της $f$ και να ανακαλύψετε το ακρότατο.
    \item Για ποιες τιμές του $a>0$, η ελάχιστη τιμή της συνάρτησης είναι ίση ή μεγαλύτερη του μηδενός;
    \item Αν $a=e$, δείξτε ότι $e^{x-1} \ge x$ για κάθε $x \in \mathbb{R}$.
    \item Χρησιμοποιώντας την προηγούμενη ανισότητα για $x= \dint_0^1 f(t) dt$, συνθέστε μια γενική ανισοτική ερμηνεία. 
\end{erwthma}
\end{thema}

%---------------------------------------------------------------
\begin{thema}{Δ}
Είναι γνωστό ότι η συνάρτηση $f$ ικανοποιεί τη σχέση $e^{f(x)} + f(x) = x$ για κάθε $x \in \mathbb{R}$. 
\begin{erwthma}
    \item Δείξτε ότι η συνάρτηση εξασφαλίζει μοναδικότητα $1-1$ και καθορίστε τον τύπο της αντίστροφης $f^{-1}(y)$.
    \item Με τη βοήθεια της παραγώγου $f'$ να μελετήσετε τη μονοτονία. (Χρησιμοποιώντας διαφοροποίηση της αρχικής σχέσης). 
    \item Να εντοπίσετε το $\lim_{x\to+\infty} f(x)$ αν γνωρίζετε ότι συγκλίνει, ή αποκλίνει και πού.
    \item Να υπολογίσετε το εμβαδόν μεταξύ της καμπύλης $C_{f^{-1}}$, του άξονα των $x$ και των ευθειών συγκράτησης $x=1, x=e+1$.
\end{erwthma}
\end{thema}

%---------------------------------------------------------------
\begin{thema}{Δ}
Δίνεται το ολοκλήρωμα $I_\nu = \dint_0^1 x^\nu e^x dx, \quad \nu \in \mathbb{N}^*$.
\begin{erwthma}
    \item Να υπολογίσετε το $I_1$ και το $I_2$.
    \item Να αποδείξετε την αναδρομική σχέση: $I_\nu = e - \nu I_{\nu-1}$.
    \item Χωρίς να υπολογίσετε τα ολοκληρώματα, να δείξετε ότι η ακολουθία $I_\nu$ είναι φθίνουσα. ($I_{\nu+1} < I_\nu$).
    \item Να υπολογίσετε το όριο $\displaystyle\lim_{\nu \to +\infty} I_\nu$.
\end{erwthma}
\end{thema}

%---------------------------------------------------------------
\begin{thema}{Δ}
Μια συνάρτηση $f: [0, +\infty) \to \mathbb{R}$ έχει συνεχή παράγωγο και ισχύει $f(0)=1$, καθώς και:
\[ f'(x) = \frac{1}{1 + f^2(x)} \quad \text{για κάθε } x \ge 0 \]
\begin{erwthma}
    \item Να αποδείξετε ότι η $f$ είναι γνησίως αύξουσα. 
    \item Ολοκληρώνοντας τη σχέση $(1+f^2(x))f'(x) = 1$, βρείτε τη σχέση που συνδέει τα $f(x)$ και $x$.
    \item Αν $x>0$, να αποδείξετε ότι ισχύει η διπλή ανισότητα: $\frac{x}{1+f^2(x)} < f(x)-1 < x$.
    \item Να λύσετε την εξίσωση $f(x) + f^3(x) = x + 2$ βασιζόμενοι στο (Δ2).
\end{erwthma}
\end{thema}

%---------------------------------------------------------------
\begin{thema}{Δ}
Δίνεται η συνάρτηση $f(x) = \sqrt{x^2 - x + 1} - a x$, με $a \in \mathbb{R}$.
\begin{erwthma}
    \item Αν $\lim_{x\to+\infty} f(x) = \frac{1}{2}$, βρείτε την παράμετρο $a$. \textit{($a=1$)}.
    \item Για $a=1$, να προσδιορίσετε την οριζόντια ασύμπτωτη της $C_f$ στο $+\infty$. 
    \item Για $a=1$, να βρείτε την πλάγια ασύμπτωτη στο $-\infty$.
    \item Υπολογίστε το $\dint_0^1 \frac{2x - 1}{2\sqrt{x^2-x+1}} dx$.
\end{erwthma}
\end{thema}

%---------------------------------------------------------------
\begin{thema}{Δ}
Έστω η $f(x) = x \sin x + \cos x$, στο $[0, 2\pi]$.
\begin{erwthma}
    \item Βρείτε την $f'(x)$ και προσδιορίστε τα διαστήματα μονοτονίας της $f$ στο $[0, 2\pi]$.
    \item Αποδείξτε ότι η $f$ παρουσιάζει ένα τοπικό μέγιστο και ένα τοπικό ελάχιστο στο \textbf{εσωτερικό} του διαστήματος $[0, 2\pi]$.
    \item Δείξτε ότι για κάθε $c \in \left(f(\pi), f\left(\frac{\pi}{2}\right)\right)$, η εξίσωση $f(x) = c$ έχει τουλάχιστον μία ρίζα στο $(0, 2\pi)$. 
    \item Να υπολογιστεί το εμβαδόν μεταξύ του $C_f$ και του $C_{f'}$ όπου χρειάζεται.
\end{erwthma}
\end{thema}

%---------------------------------------------------------------
\begin{thema}{Δ}
\textbf{Εναλλακτικό Πρόβλημα Ρυθμού Μεταβολής \& Ανάλυσης}\\
Ένα σημείο $M$ κινείται πάνω στον κλάδο της υπερβολής $y = \frac{1}{x}$, με $x>0$. Η τετμημένη $x$ του $M$ αυξάνεται με σταθερό ρυθμό $2 \text{ cm/s}$ ($x'(t) = 2$). 
\begin{erwthma}
    \item Να βρείτε το ρυθμό μεταβολής της τεταγμένης $y(t)$ τη χρονική στιγμή $t_0$ που το $M$ διέρχεται από το σημείο $\left(2, \frac{1}{2}\right)$.
    \item Θεωρούμε το τρίγωνο $OAM$ όπου $O(0,0)$ και $A$ η προβολή του $M$ στον $x'x$. Ποιος ο ρυθμός μεταβολής του εμβαδού του $OAM$;
    \item Η εφαπτομένη ($\epsilon$) στο $M$ τέμνει τους άξονες στα σημεία $\Gamma, \Delta$. Να δείξετε ότι το εμβαδόν του σχηματιζόμενου τριγώνου $O\Gamma\Delta$ είναι σταθερό (ανεξάρτητο του χρόνου $t$) και ισούται με $2$.
    \item Να εξετάσετε τη μονοτονία της απόστασης $(OM)$ ως συνάρτηση του $x$. Σε ποιο σημείο πλησιάζει εγγύτερα το $M$ την αρχή των αξόνων;
\end{erwthma}
\end{thema}

\begin{thema}{Δ}
Έστω συνάρτηση $f:\mathbb{R}\to\mathbb{R}$ για την οποία ισχύει $f^3(x) + f(x) = x$ για κάθε $x \in \mathbb{R}$.
\begin{erwthma}
    \item Να αποδείξετε ότι η συνάρτηση $f$ αντιστρέφεται στο $\mathbb{R}$ και να προσδιορίσετε τον τύπο της αντίστροφης $f^{-1}$.
    \item Να αποδείξετε ότι η $f$ είναι γνησίως αύξουσα και να προσδιορίσετε το μοναδικό σημείο τομής της γραφικής της παράστασης $C_f$ με τους άξονες.
    \item Να αποδείξετε ότι η $f$ είναι συνεχής στο $\mathbb{R}$. (Μπορείτε να μελετήσετε διαφορές της μορφής $f(x)-f(x_0)$ σε συνάρτηση με το $x-x_0$).
    \item Να υπολογίσετε το $\displaystyle\lim_{x\to +\infty} \frac{f(x)}{x}$ και να μελετήσετε αν η $C_f$ δέχεται πλάγια ασύμπτωτη.
    \item Να υπολογίσετε το εμβαδόν του χωρίου που περικλείεται από την αντίστροφη συνάρτηση $C_{f^{-1}}$, τον άξονα $x'x$, και την ευθεία $x=1$. Στη συνέχεια να κυκλώσετε το άθροισμα $\dint_0^{1} f(x)dx + \dint_{f(0)}^{f(1)} f^{-1}(x)dx$.
\end{erwthma}
\end{thema}

%---------------------------------------------------------------
\begin{thema}{Δ}
Θεωρούμε μια δύο φορές παραγωγίσιμη συνάρτηση $f: (0, +\infty) \to \mathbb{R}$ που ικανοποιεί τις συνθήκες:
\[ x^2 f''(x) + 2x f'(x) = 1 \quad \text{για κάθε } x > 0, \quad f(1)=1 \text{ και } f'(1)=-1 \]
\begin{erwthma}
    \item Να προσδιορίσετε τον τύπο της συνάρτησης, αποδεικνύοντας ότι $f(x) = \ln x - 2x + 3$.
    \item Να αποδείξετε ότι η εξίσωση $f(x) = 0$ έχει ακριβώς δύο ρίζες, μία στο διάστημα $(0,1)$ και μία στο $(1, +\infty)$.
    \item Έστω $\rho_1, \rho_2$ οι ρίζες του ερωτήματος $\Delta 2$ με $\rho_1 < 1 < \rho_2$. Να δείξετε ότι η $f(x)$ είναι κοίλη και να προσδιορίσετε την απάντηση για το πρόσημό της στο $[\rho_1, \rho_2]$.
    \item Να αποδείξετε τη σχέση $\displaystyle \dint_{\rho_1}^{\rho_2} f(x) dx < \frac{1}{2}(\rho_2 - \rho_1)(\rho_2 + \rho_1 - 4)$.
\end{erwthma}
\end{thema}

%---------------------------------------------------------------
\begin{thema}{Δ}
Δίνεται η συνάρτηση $f(x) = e^x + x - 1$, με $x \in \mathbb{R}$.
\begin{erwthma}
    \item Να αποδείξετε ότι η εξίσωση $f(x)=0$ έχει μοναδική ρίζα την $x=0$, και στη συνέχεια να λύσετε την ανίσωση $f(f(e^x - 1)) > f(e - 1)$.
    \item Να υπολογίσετε το $\displaystyle\lim_{x\to +\infty} \frac{x f(x) + \ln x}{f(x) + e^{x-1}}$.
    \item Να αποδείξετε ότι για κάθε $x>0$ ισχύει: $x e^x > e^x - 1$. 
    \item Να λύσετε στο $(0, +\infty)$ την εξίσωση $e^x(x^2 - x) + x = \ln\left(\frac{e^x - 1}{x}\right)$.
\end{erwthma}
\end{thema}

%---------------------------------------------------------------
\begin{thema}{Δ}
Θεωρούμε την συνάρτηση $f(x) = \frac{\ln x}{x} - \frac{1}{x} + a$, $x>0, a \in \mathbb{R}$.
\begin{erwthma}
    \item Αν γνωρίζετε ότι η γραφική παράσταση της $f'$ τέμνει τον άξονα $x$ στο σημείο $x_0 = e^{3/2}$, να αποδείξετε ότι, με $a=0$, η $f$ παρουσιάζει τοπικό ακρότατο (το ποιο θα το ανακαλύψετε) στο σημείο αυτό.
    \item Να αποδείξετε ότι η συνάρτηση $g(x) = x^3 f(x)$ είναι γνησίως αύξουσα στο $[1, e]$.
    \item Αν ισχύει η αρχική υπόθεση με $a=1$, να αποδείξετε ότι υπάρχει μοναδικό $\rho \in (1, e)$ τέτοιο ώστε η $f$ να μηδενίζεται.
    \item Να αποδείξετε ότι $\dint_\rho^e f(x) dx > 0$, με $\rho$ το σημείο του $\Delta3$.
\end{erwthma}
\end{thema}

%---------------------------------------------------------------
\begin{thema}{Δ}
Έστω η παραγωγίσιμη συνάρτηση $f:(0, +\infty) \to \mathbb{R}$ με $f(e)=1$ τέτοια ώστε $x \ln x \cdot f'(x) + f(x) = \ln x$ για κάθε $x \in (1, +\infty)$.
\begin{erwthma}
    \item Να αποδείξετε ότι $f(x) = \frac{x\ln x - x + e}{\ln x}$ για κάθε $x > 1$.
    \item Να δείξετε ότι η $f$ είναι γνησίως φθίνουσα στο $(1, +\infty)$.
    \item Να προσδιορίσετε την οριζόντια ασύμπτωτη της $C_f$ στο $+\infty$ καθώς και την κατακόρυφη μορφή της όταν $x \to 1^+$.
    \item Θεωρούμε την εφαπτομένη της $C_f$ στο $\xi = e$. Δείξτε ότι η εφαπτομένη αυτή χωρίζει σε δύο ίσα μέρη το εμβαδόν ενός τριγώνου, το οποίο θα προσδιορίσετε βρίσκοντας τις τομές της με τους άξονες $Oxy$.
\end{erwthma}
\end{thema}

%---------------------------------------------------------------
\begin{thema}{Δ}
Δίνονται οι συναρτήσεις $f(x) = e^x - \alpha x$ και $g(x) = \ln x + \alpha x$ με $\alpha > 0$. Είναι γνωστό ότι η $f$ παρουσιάζει ελάχιστο ίσο με $0$, ενώ η $g$ έχει κοινή εφαπτομένη με την $f$ σε ένα σημείο $x_0$.
\begin{erwthma}
    \item Να αποδείξετε την τιμή του $\alpha = e$.
    \item Να δείξετε ότι η εξίσωση $e x^2 + 1 = e^x$ έχει ως μοναδικές λύσεις το $0$ και το $1$. 
    \item Να μελετήσετε την $g(x)$ ως προς την κοιλότητα και να εξηγήσετε γιατί οι καμπύλες δεν μπορούν να τέμνονται σε περισσότερα από δύο σημεία.
    \item Να υπολογίσετε το εμβαδόν μεταξύ της $f(x)$ και του άξονα αν αντικαταστήσουμε το $a=e$ στη σχετική τομή του Δ2.
\end{erwthma}
\end{thema}

%---------------------------------------------------------------
\begin{thema}{Δ}
H συνεχής συνάρτηση $f:\mathbb{R}\to\mathbb{R}$ παρουσιάζει την ιδιότητα: $(f(x) - x)(f(x) + x) = x^2 + 1$ για κάθε $x \in \mathbb{R}$ και διέρχεται από το σημείο $(0, 1)$.
\begin{erwthma}
    \item Να δείξετε ότι ο τύπος της συνάρτησης είναι $f(x) = \sqrt{2x^2 + 1}$.
    \item Να δείξετε ότι υπάρχουν $\xi_1, \xi_2 \in (-1, 1)$ τέτοια ώστε $f'(\xi_1) + f'(\xi_2) = 0$.
    \item Να δείξετε ότι τα κοίλα της $f$ στρέφονται προς τα άνω (strict Convex).
    \item Αν $\alpha < \beta < \gamma$ πραγματικοί αριθμοί σε αριθμητική πρόοδο, δείξτε ότι $f(\alpha) + f(\gamma) > 2f(\beta)$.
\end{erwthma}
\end{thema}

%---------------------------------------------------------------
\begin{thema}{Δ}
Δίνεται η συνάρτηση $f(x) = \begin{cases} \frac{x^2 \sin(1/x)}{x}, & x \ne 0 \\ 0, & x = 0 \end{cases}$. (Προσοχή στον ορισμό σε μηδενικά).
\begin{erwthma}
    \item Να δείξετε ότι η $f$ είναι συνεχής στο $0$.
    \item Να αποδείξετε τη δυνατότητα παραγώγισης στο $x_0 = 0$ και να υπολογίσετε την $f'(0)$.
    \item Είναι η $f'$ συνεχής στο 0; Να δικαιολογήσετε την απάντηση.
    \item Να προσδιορίσετε αν η $f$ παρουσιάζει ακρότατο στο $0$. (Αυστηρή τεκμηρίωση με εναλλαγές προσήμου παραγώγου κοντά στο μηδέν).
\end{erwthma}
\end{thema}

%---------------------------------------------------------------
\begin{thema}{Δ}
Δίνεται $f: (0,+\infty) \to \mathbb{R}$ παραγωγίσιμη με $f(1)=1$ και $x f'(x) - f(x) = x^2 e^{-x}$ για κάθε $x>0$.
\begin{erwthma}
    \item Να αποδείξετε ότι $f(x) = x (1 + e^{-1} - e^{-x})$.
    \item Να υπολογίσετε το $\displaystyle\lim_{x\to +\infty} \frac{f(x)}{x}$. 
    \item Δείξτε ότι η $f$ έχει μία και μοναδική πλάγια ασύμπτωτη στο $+\infty$.
    \item Να βρείτε την εξίσωση της εφαπτομένης της $C_f$ η οποία είναι παράλληλη της πλάγιας ασύμπτωτης.
    \item Να υπολογίσετε το ολοκλήρωμα $\dint_1^2 \frac{f(x)}{x} dx$.
\end{erwthma}
\end{thema}

%---------------------------------------------------------------
\begin{thema}{Δ}
Έστω συνάρτηση $f:\mathbb{R} \to \mathbb{R}$ συνεχής, η οποία ικανοποιεί τη σχέση $\dint_0^{x} f(t) dt + f(x) = x e^x$ για όλους τους πραγματικούς.
\begin{erwthma}
    \item Παραγωγίζοντας κατάλληλα τα μέλη της ισότητας (εξηγήστε γιατί η $f$ οφείλει να είναι παραγωγίσιμη), αποδείξτε ότι  $f'(x) + f(x) = e^x(x+1)$.
    \item Λύνοντας τη διαφορική της παραπάνω σχέσης, να δείξετε ότι $f(x) = x e^x$.
    \item Να λύσετε την ανίσωση $x e^x > \ln(x+1)(x+1)$ για $x>0$. 
    \item Να δείξετε ότι το εμβαδόν $E(a)$ μεταξύ της καμπύλης και του άξονα των $x$ στο διάστημα $[-a, 0]$ τείνει στο $1$ όταν $a \to +\infty$.
\end{erwthma}
\end{thema}

%---------------------------------------------------------------
\begin{thema}{Δ}
Μία συνάρτηση $f:(0, +\infty) \to \mathbb{R}$ είναι δύο φορές παραγωγίσιμη, ικανοποιώντας την $f(x) + f''(x) = 0$ και τις οριακές τιμές $f(\pi) = 0, \ f'(\pi) = -1$.
\begin{erwthma}
    \item Αποδείξτε (θεωρώντας τη συνάρτηση $g(x) = (f(x) - \sin x)^2 + (f'(x) - \cos x)^2$) ότι $f(x) = \sin x$.
    \item Να αποδείξετε ότι η εξίσωση $f(x) = x - 2$ έχει μοναδική ρίζα $\rho \in (0, \pi)$.
    \item Να υπολογίσετε την οριακή συμπεριφορά $\displaystyle\lim_{x\to \rho} \frac{f(x) - x + 2}{\sin(x - \rho)}$.
    \item Αν $\Omega$ είναι το χωρίο $x \in [0, \pi]$, να βρείτε πόσο αλλάζει το εμβαδόν αν τη συνάρτηση $f$ αντικαταστήσουμε από την $\left|f'(x)\right|$.
\end{erwthma}
\end{thema}

%---------------------------------------------------------------
\begin{thema}{Δ}
Η συνάρτηση $f$ ορίζεται ως $f(x) = e^{x-1} + \ln(x) + x - 2$ με $x>0$.
\begin{erwthma}
    \item Να εξετάσετε τη μονοτονία της και να δείξετε ότι είναι αντίστρεψιμη (1-1). 
    \item Να λύσετε την εξίσωση $e^{\ln x - 1} + \ln(\ln x) + \ln x = 2$.
    \item Να υπολογίσετε την παράγωγο της $f^{-1}$ στο σημείο $(f(1))$, με δεδομένο ότι η $f$ είναι παραγωγίσιμη στο πεδίο ορισμού και η αντίστροφη είναι ομοίως παραγωγίσιμη.
    \item Δείξτε ότι υπάρχει ακριβώς ένα $x_0 \in (1, e)$ τέτοιο ώστε: $f(x_0) = \frac{e^2 - 1}{x_0}$.
\end{erwthma}
\end{thema}

%---------------------------------------------------------------
\begin{thema}{Δ}
Ένας σφαίρος πάγου αρχίζει να λιώνει. Αν ο ρυθμός μείωσης του όγκου του είναι ανάλογος της επιφάνειάς του. Θεωρούμε την ακτίνα $R(t)$ ως συνάρτηση του χρόνου.
\begin{erwthma}
    \item Να δείξετε ότι ο ρυθμός μείωσης της ακτίνας είναι σταθερός διαχρονικά ($R'(t) = -k$).
    \item Αν αρχικά ο όγκος είναι $V_0$ και μετά από $\Delta t = 2$ ώρες είναι $V_0 / 8$, ποιος είναι ο χρόνος ολικής τήξης;
    \item Να μοντελοποιήσετε συνάρτηση $f(x) = e^{k x}$ και να δείξετε ότι ο λόγος $\frac{f(x+c)}{f(x)}$ παραμένει σταθερός ανεξάρτητα του $x$. (Ασύγκριτος συνδυασμός ρυθμού με ιδιότητα εκθετικής).
    \item Να χρησιμοποιήσετε το Θεώρημα Bolzano για να λύσετε την εξίσωση του υπολοίπου του όγκου ίσω με $1$ αν $V_0 = 8\pi$.
\end{erwthma}
\end{thema}

%---------------------------------------------------------------
\begin{thema}{Δ}
Δίνεται η εξίσωση $(x^2+1) f''(x) + x f'(x) - f(x) = 0$ με αρχικές συνθήκες $f(0)=0, f'(0)=1$ σε όλο το $\mathbb{R}$.
\begin{erwthma}
    \item Υποθέτοντας ότι η $f(x)$ είναι πολυώνυμο βαθμού 1, βρείτε έστω μία λύση $f(x) = c x + d$. Ανταποκρίνεται $f(x)=x$;
    \item Αποδείξτε ότι η συνάρτηση διατηρεί συγκεκριμένη κυρτότητα στο $(0, +\infty)$.
    \item Λύστε την ανίσωση $f(e^x - 1) > f(x)$.
    \item Υπολογίστε το ολοκλήρωμα $\dint_0^1 x^2 f(x) dx$ γνωρίζοντας ότι $f(x)=x$.
\end{erwthma}
\end{thema}

%---------------------------------------------------------------
\begin{thema}{Δ}
Έστω ότι $g: (0, +\infty) \to \mathbb{R}$ με $g(x) = \frac{\ln x + \ln(1/x)}{x^2}$. Και έστω $f(x)$ μια άλλη τέτοια που $f'(x) = g(x)$.
\begin{erwthma}
    \item Παρατηρήστε την ιδιότητα του λογαρίθμου για να δείξετε ότι $f'(x) = 0$. Άρα η $f$ είναι σταθερή; Ναι ή όχι; 
    \item Αν $f(1)= \int_0^1 e^{-t^2} dt + c$, τι εκπροσωπεί αυτό; 
    \item Προσομοιώστε πλέον τη συνάρτηση $h(x) = \ln(x) + x$. Δείξτε ότι $\dint_1^e h(x) dx > 2$.
    \item Αποδείξτε ότι δεν υπάρχει εφαπτομένη της $h(x)$ διερχόμενη από το $(0, -\pi)$.
\end{erwthma}
\end{thema}

\begin{thema}{Δ}
Δίνεται η παραγωγίσιμη συνάρτηση $f : \mathbb{R} \to \mathbb{R}$ για την οποία ισχύει $e^{f(x)} + f(x) = x + 1$ για κάθε $x \in \mathbb{R}$.
\begin{erwthma}
    \item Αποδείξτε ότι η $f$ είναι γνησίως αύξουσα στο $\mathbb{R}$. Ποιο είναι το σύνολο τιμών της;
    \item Αποδείξτε ότι $f(x) < x$ για κάθε $x > 0$.
    \item Να βρείτε, εάν υπάρχει, το $\displaystyle\lim_{x\to +\infty} f(x)$.
    \item Να λύσετε την ανίσωση $\displaystyle e^{f(x^2 - 4)} - e^{f(3x)} < f(3x) - f(x^2 - 4)$.
    \item Έστω $E$ το εμβαδόν του χωρίου μεταξύ $C_f$, του άξονα των $x$ και των ευθειών $x=0, x=e+1$. Αποδείξτε ότι $E = \frac{e}{2}$.
\end{erwthma}
\end{thema}

%---------------------------------------------------------------
\begin{thema}{Δ}
Έστω ότι η $g: (0, +\infty) \to \mathbb{R}$ ικανοποιεί τη σχέση $g(xy) = g(x) + g(y)$ για όλα τα $x, y > 0$. Δίνεται επιπλέον ότι η $g$ είναι παραγωγίσιμη στο $x_0 = 1$ με $g'(1) = 1$.
\begin{erwthma}
    \item Να αποδείξετε ότι η $g$ είναι παραγωγίσιμη σε κάθε $x_0 > 0$ και να δείξετε ότι $g'(x) = \frac{1}{x}$.
    \item Να δείξετε ότι $g(x) = \ln x$. 
    \item Έστω $f(x) = g(x) - x + 1$. Δείξτε ότι $f(x) \le 0$ και λύστε την εξίσωση $x^{e} = e^x$.
    \item Να υπολογίσετε το $\displaystyle\lim_{x\to 0^+} x g(x)$.
\end{erwthma}
\end{thema}

%---------------------------------------------------------------
\begin{thema}{Δ}
Η συνάρτηση $h(x) = a^x - x^a$ ($a > 0, a \ne 1$) ορίζεται στο $(0, +\infty)$. Αν η εξίσωση $h(x) = 0$ έχει διπλή ρίζα στο $(0, +\infty)$:
\begin{erwthma}
    \item Να αποδείξετε ότι ο πραγματικός αριθμός $a$ ισούται με $e$.
    \item Αντικαθιστώντας $a=e$, σχεδιάστε τον πίνακα μεταβολών της $\phi(x) = e^x - x^e$ και προσδιορίστε τα ακρότατα.
    \item Αποδείξτε ότι $\phi(x) > 0$ για $x \in (0, e) \cup (e, +\infty)$.
    \item Να υπολογιστεί το εμβαδόν μεταξύ των $y=e^x, y=x^e$ στο $[1, e]$.
\end{erwthma}
\end{thema}

%---------------------------------------------------------------
\begin{thema}{Δ}
Έστω η παραγωγίσιμη συνάρτηση $f: \mathbb{R} \to \mathbb{R}$ η οποία ικανοποιεί την κατάσταση:
$(x^2 + 1) f'(x) + 2x f(x) = 2x$ με $f(0) = 0$.
\begin{erwthma}
    \item Παρατηρώντας το αριστερό μέλος ως παράγωγο γινομένου, λύστε τη διαφορική εξίσωση για να αποδείξετε ότι $f(x) = \frac{x^2}{x^2+1}$.
    \item Αποδείξτε (εφαρμόζοντας διαδοχικά θεωρήματα στην $f$) ότι υπάρχει μοναδικό σημείο καμπής στον θετικό ημιάξονα το οποίο θα προσδιορίσετε.
    \item Δείξτε ότι υπάρχει $x_0 \in (0, 1)$ ώστε η εφαπτομένη να είναι παράλληλη με τη χορδή $AB$ όπου $A(0,0)$ και $B(1, 1/2)$.
    \item Χωρίς να ολοκληρώσετε, προσδιορίστε αναλυτικά ένα κάτω φράγμα (μία ελάχιστη τιμή) για το $\dint_0^1 f(x)dx$.
\end{erwthma}
\end{thema}

%---------------------------------------------------------------
\begin{thema}{Δ}
Θεωρούμε τη συνάρτηση $f(x) = x \ln x$. Δίνεται επίσης μια δεύτερη συνάρτηση $g$ παραγωγίσιμη ώστε $g'(x) = f(x)$.
\begin{erwthma}
    \item Να μελετήσετε την μονοτονία της $f$ και να δείξετε ότι $f(x) \ge -1/e$.
    \item Να δείξετε ότι η συνάρτηση $g$ παρουσιάζει ένα τοπικό ελάχιστο στο $x=1$.
    \item Υπολογίστε τον ακριβή τύπο της $g(x)$, αν περνά από το σημείο $(1, -1/4)$.
    \item Δείξτε ότι $\dint_1^e g(x) dx > 0$.
\end{erwthma}
\end{thema}

%---------------------------------------------------------------
\begin{thema}{Δ}
Δίνονται οι συναρτήσεις $f(x) = e^x - 1$ και $g(x) = -\frac{1}{x+1}$.
\begin{erwthma}
    \item Αποδείξτε ότι οι $C_f, C_g$ τέμνονται σε ένα μόνο σημείο το οποίο και να προσδιορίσετε.
    \item Δείξτε ότι υπάρχει κοινή εφαπτομένη τους σε ένα σημείο και βρείτε την εξίσωσή της.
    \item Αποδείξτε ότι $f(x) \ge g(x)$ για $x \ge 0$.
    \item Υπολογίστε το εμβαδόν μεταξύ της $f(x)$ και της $g(x)$ από $x=0$ μέχρι $x=1$.
\end{erwthma}
\end{thema}

%---------------------------------------------------------------
\begin{thema}{Δ}
Έστω η πολυωνυμική συνάρτηση $P(x) = x^3 - 3x^2 + \lambda x + \mu$. 
\begin{erwthma}
    \item Να βρείτε σχέση των $\lambda, \mu$ ώστε το $x=1$ να είναι σημείο καμπής.
    \item Αν $\lambda=3, \mu=-1$, μελετήστε το $P(x)$ για ρίζες. Πόσες πραγματικές ρίζες έχει;
    \item Δείξτε ότι το ολοκλήρωμα $\dint_0^2 \left(P(x)\right)^{2023} dx = 0$. (Ιδιότητα συμμετρίας σε εκφράσεις της μορφής $\phi(a-x)$).
    \item Προσδιορίστε τα όρια της ταχύτητας όταν $P'(x) \to +\infty$.
\end{erwthma}
\end{thema}

%---------------------------------------------------------------
\begin{thema}{Δ}
Δίνεται η $f(x) = x + \sin x$. 
\begin{erwthma}
    \item Αποδείξτε ότι η $f$ είναι γνησίως αύξουσα. 
    \item Να λύσετε την ανίσωση $x + \sin x > \pi$.
    \item Να δείξετε ότι τα $f''(x), f^{(4)}(x)$ εναλλάσσουν πρόσημα σε συγκεκριμένες θέσεις, ορίζοντας άπειρα σημεία καμπής.
    \item Να υπολογίσετε το εμβαδόν μεταξύ $f(x), x=0, x=\pi, y=x$.
\end{erwthma}
\end{thema}

%---------------------------------------------------------------
\begin{thema}{Δ}
Έστω συνεχής παραγωγίσιμη συνάρτηση $f : [-1, 1] \to \mathbb{R}$ όπου $f^2(x) + x^2 = 1$ για κάθε $x \in [-1, 1]$.
\begin{erwthma}
    \item Εάν $f(0) = 1$, αποδείξτε αναλυτικά ότι $f(x) = \sqrt{1-x^2}$ για κάθε $x \in [-1, 1]$.
    \item Εντοπίστε τις ασύμπτωτες (αν υπάρχουν) της $g(x) = \frac{1}{f(x)}$ στο $(-1, 1)$. 
    \item Προσδιορίστε το εμβαδόν του κυκλικού τομέα που ορίζει το $f(x)$ με τον άξονα $x$ και τις ευθείες $x=0, x=1$.
    \item Λύστε την ανίσωση $f(x) < x$. 
\end{erwthma}
\end{thema}

%---------------------------------------------------------------
\begin{thema}{Δ}
Συνάρτηση $f: \mathbb{R} \to \mathbb{R}$ ικανοποιεί την: $e f(x) + e^{-f(x)} = 2x$.
\begin{erwthma}
    \item Βρείτε το πεδίο τιμών, και αν δεχτούμε ότι είναι 1-1, βρείτε την αντίστροφή της $\phi(y)$.
    \item Βρείτε το πεδίο ορισμού της αντίστροφης και καθορίστε τις ιδιότητές της (μονοτονία, κυρτότητα).
    \item Αν $\phi(y)$ είναι γνωστή, υπολογίστε τα όρια $\lim_{y \to +\infty} \frac{\phi(y)}{e^y}$.
    \item Σε ένα υπολογιστικό χωρίο συνδεδεμένο με το $\phi(y)$, υπολογίστε το ολοκλήρωμα $\dint_{\sqrt{e}}^{e} \phi(x) dx$.
\end{erwthma}
\end{thema}

%---------------------------------------------------------------
\begin{thema}{Δ}
Θεωρούμε την οικογένεια συναρτήσεων $f_\lambda(x) = \frac{e^x}{x^\lambda}$ με $x>0, \lambda>0$.
\begin{erwthma}
    \item Μελετήστε τα ακρότατα της $f_\lambda(x)$ συναρτήσει του $\lambda$. Σε ποιο σημείο μεγιστοποιείται/ελαχιστοποιείται;
    \item Προσδιορίστε τη μέγιστη κάτω ρίζα του διαφορικού.
    \item Επιλέξτε $\lambda = 1$. Υπολογίστε το όριο $\lim_{x \to +\infty} f_1(x)$.
    \item Επιφυλάσσοντας το $\lambda = e$, λύστε την ανίσωση $f_e(x) \ge e^e$.
\end{erwthma}
\end{thema}

%---------------------------------------------------------------
\begin{thema}{Δ}
Ο ρυθμός αύξησης πληθυσμού βακτηρίων $P(t)$ δίνεται από τον νόμο $P'(t) = 0.5 P(t)$, και αρχικά έχουμε $P(0) = 1000$.
\begin{erwthma}
    \item Προσδιορίστε τη συνάρτηση πληθυσμού.
    \item Ένα φάρμακο αρχίζει να εφαρμόζεται τη στιγμή $t=10$ ανατρέποντας το ρυθμό σε $P'(t) = -0.5 P(t)$. Βρείτε τον νέο τύπο πληθυσμού μετά από $t=10$.
    \item Σε ποια χρονική στιγμή ο πληθυσμός θα επιστρέψει στα 1000 βακτήρια;
    \item Αυτό το φαινόμενο της συνεχειοποίησης σε δύο κλάδους πώς γεωμετρικά ερμηνεύεται ως προς την $C_P$; Είναι παραγωγίσιμη στο $t=10$;
\end{erwthma}
\end{thema}

%---------------------------------------------------------------
\begin{thema}{Δ}
Οι συναρτήσεις $f(x), g(x)$ έχουν κοινό πεδίο ορισμού και $f'(x) = g(x)$ και $g'(x) = f(x)$, με $f(0)=1, g(0)=0$.
\begin{erwthma}
    \item Δείξτε ότι $(f(x)-g(x))'=-(f(x)-g(x))$ και $(f(x)+g(x))'= f(x)+g(x)$.
    \item Βρείτε τους τύπους των $f$ και $g$.
    \item Αποδείξτε ότι $\forall x \in \mathbb{R}$, $f^2(x) - g^2(x) = 1$.
    \item Σχεδιάστε μια διαδικασία ολοκλήρωσης βρίσκοντας το $\dint_0^1 f^2(x) dx$.
\end{erwthma}
\end{thema}

%---------------------------------------------------------------
\begin{thema}{Δ}
Έστω $f: \mathbb{R} \to \mathbb{R}$ συνεχής και $F(x)$ μια αρχική (παράγουσα) της $f(x)$. Αν ισχύει $F(x) \cdot f(x) = x$ και $F(0)=1$.
\begin{erwthma}
    \item Παραγωγίζοντας τη σχέση $(1/2 F^2(x))' = x$, δείξτε ότι $F(x) = \sqrt{x^2+1}$.
    \item Εκφράστε την $f(x)$. Είναι άρτια, περιττή;
    \item Ελέγξτε τις κατακόρυφες και οριζόντιες ασύμπτωτες των $F(x), f(x)$.
    \item Υπολογίστε επιφανειακό ολοκλήρωμα $\dint_0^1 x F(x) f(x) dx$.
\end{erwthma}
\end{thema}

%---------------------------------------------------------------
\begin{thema}{Δ}
Η συνάρτηση $h(x) = x^3 \ln(x^2) - x + 1$ ελέγχεται ως προς τα κριτήρια ακροτάτων.
\begin{erwthma}
    \item Δείξτε ότι έχει πεδίο ορισμού το $\mathbb{R}^*$ και είναι άρτια; (Όχι, $\ln(-x^2)$... είναι $\ln(|x|^2)$). Βρείτε το συμμετρικό πεδίο.
    \item Να βρείτε τις τοπικές ασύμπτωτες.
    \item Η παράγωγος στο $x_0 = \sqrt{e}$ δίνεται. Υπολογίστε την εφαπτομένη.
    \item Αποδείξτε, χωρίς τον υπολογισμό, ότι η εξίσωση $h(x) = 0$ έχει τουλάχιστον μια ρίζα. Καθώς και την ενσωματωμένη έκφανσή της σε Θ.Ε.Τ.
\end{erwthma}
\end{thema}
\end{document}